\section{Factor branching framework}
\label{sec:method}

\subsection{Framework description}
\label{subsec:branching-framework}

% motivation for breaking down a system of polynomial equations
Polynomial-ideal computations can be intractable even for small systems~\cite{mayr1982complexity}. Thus, decomposing a system into smaller and simpler
branches may make the individual solver calls easier.

% How breaking down is done
We break a system into smaller branches by factoring each generator and
enumerating the minimal hitting sets of the factor sets. Each hitting set
generates a new branch, which we interreduce and factor again, until no
branch changes.

Algorithm~\ref{alg:Refine} first factors $P$ and then applies these two
operations to every branch until it reaches a fixed point. We call a branch
\emph{terminal} if one round leaves it unchanged.

\begin{algorithm}[htbp]
\caption{Refinement}
\label{alg:Refine}
\begin{algorithmic}[1]
    \Require A finite polynomial system $P$
    \Ensure The terminal branches of $P$
    \State $W\gets\mathcal H(P)$
    \Repeat
        \State $W_{\mathrm{prev}}\gets W$, \quad $W\gets\emptyset$
        \ForAll{$Q\in W_{\mathrm{prev}}$}
            \State $Q\gets\Call{Interreduce}{Q}$
            \State $W\gets W\cup\mathcal H(Q)$
        \EndFor
    \Until{$W=W_{\mathrm{prev}}$}
    \State $W\gets\Call{Prune}{W}$ \Comment{Optional}
    \State \Return $W$
\end{algorithmic}
\end{algorithm}

% All about pruning
Hitting sets are minimal only among the children of one branch, so the loop
of Algorithm~\ref{alg:Refine} can end with distinct branches $H$ and $H'$
such that $V(H)\subseteq V(H')$. Solving $H$ then adds no solutions.
$\Call{Prune}{W}$ returns the branches whose solution sets it cannot certify to be contained in each other.

% First interreduce?
Interreducing $P$ first can expose factors that its generators hide, but interreducing can also eliminate factorizations without exposing new ones.
So the decision on whether to interreduce first depends on $P$.

% How branching is done
Let $\mathcal L(P)$ denote the terminal branches returned by
Algorithm~\ref{alg:Refine} on $P$.
For each $H\in\mathcal L(P)$, the branch solver must either describe $V(H)$ exactly or certify that $V(H)=\emptyset$.

\begin{algorithm}[htbp]
\caption{Factor branching}
\label{alg:Branching}
\begin{algorithmic}[1]
    \Require A set of polynomials $P = \{p_1, \ldots, p_m\}$ and a branch solver $\mathcal{B}$.

    \State $S \gets \emptyset$ \Comment{Solution-set representations}
    \ForAll{$H\in\Call{Refine}{P}$}
        \State $S \gets S \cup \Call{$\mathcal{B}$}{H}$
    \EndFor

    \State \Return $S$
\end{algorithmic}
\end{algorithm}

% A flow chart of the method
\begin{figure}[pos=htbp]
    \centering
    \resizebox{0.72\textwidth}{!}{\begin{tikzpicture}[
    font=\small,
    every node/.style={align=center},
    blue box/.style={draw=blue!70!cyan, line width=1pt, rounded corners=6pt,
        fill=blue!50!cyan!4, inner sep=7pt},
    blue subbox/.style={draw=blue!60!cyan, line width=.6pt, rounded corners=3pt,
        fill=blue!50!cyan!10, inner sep=4pt, font=\footnotesize},
    loop box/.style={draw=orange!85!black, line width=1pt, rounded corners=6pt,
        fill=yellow!40!orange!12, inner sep=7pt},
    decision/.style={draw=green!55!black, line width=1pt, diamond, aspect=1.8,
        fill=green!60!teal!8, inner sep=1pt},
    green box/.style={draw=green!55!black, line width=1pt, rounded corners=6pt,
        fill=green!60!teal!4, inner sep=7pt},
    green subbox/.style={draw=green!55!black, line width=.6pt, rounded corners=3pt,
        fill=white, inner sep=4pt, font=\footnotesize},
    red box/.style={draw=red!70, line width=1pt, rounded corners=6pt,
        fill=red!4, inner sep=7pt},
    red subbox/.style={draw=red!70, line width=.6pt, rounded corners=3pt,
        fill=white, inner sep=4pt, font=\footnotesize},
    gray box/.style={draw=black!55, line width=1pt, rounded corners=6pt,
        fill=black!4, inner sep=7pt},
    gray subbox/.style={draw=black!55, line width=.6pt, rounded corners=3pt,
        fill=white, inner sep=4pt, font=\footnotesize},
    return box/.style={draw=violet, line width=1pt, rounded corners=6pt,
        fill=violet!6, inner sep=7pt},
    arr/.style={-{Stealth[length=2.6mm]}, line width=1pt},
]

% ---- top row: input -> preprocessing -> initialization ----
\node[blue box] (input) {%
    \textcolor{blue}{\bfseries Input}\\[3pt]
    Polynomials $P = \{p_1,\ldots,p_m\}$};

\node[blue box, right=12mm of input] (prep) {%
    \textcolor{blue}{\bfseries Preprocessing}\\[4pt]
    \tikz{\node[blue subbox] {Interreduce the polynomials};}
    \textcolor{blue}{$\boldsymbol{\rightarrow}$}
    \tikz{\node[blue subbox] {Factor into irreducible factors $C_1,\ldots,C_m$};}};

\node[blue box, right=12mm of prep] (init) {%
    \textcolor{blue}{\bfseries Initialization}\\[3pt]
    Unsatisfiable cores $C \leftarrow$ empty set\\
    Solutions $S \leftarrow$ empty set};

% ---- loop over minimal hitting sets ----
\node[loop box] (loop) at ($(init.south)+(-20mm,-16mm)$) {%
    \textcolor{orange!75!black}{\bfseries Loop}\\[3pt]
    Iterate over minimal hitting sets $HS$ of $C_1,\ldots,C_m$};

% ---- unsat-core decision + prune ----
\node[decision, below=9mm of loop] (learned) {%
    Does $HS$ contain a\\ learned unsat core\\ from $C$?};

\node[green box, right=14mm of learned] (prune) {Prune / skip};

% ---- solve ----
\node[blue box, below=9mm of learned] (solve) {%
    \textcolor{blue}{\bfseries Check satisfiability of $HS$}\\[3pt]
    \tikz{\node[blue subbox] {Gr\"obner basis computation};}};

% ---- two outcome branches ----
\node[red box] (unsat) at ($(solve.south)+(-34mm,-20mm)$) {%
    \textcolor{red}{\bfseries Unsatisfiable}\\[3pt]
    \tikz{\node[red subbox] {$U \leftarrow$ minimal unsat core of $HS$\\
        by iterative refinement,\\ learn it: $C \leftarrow C \cup \{U\}$};}};

\node[green box] (sat) at ($(solve.south)+(34mm,-20mm)$) {%
    \textcolor{green!45!black}{\bfseries Satisfiable}\\[3pt]
    \tikz{\node[green subbox] {Add the solution to $S$};}};

% ---- bottom rail ----
\coordinate (raill) at ($(unsat.south)+(0,-8mm)$);
\coordinate (railm) at (raill -| solve.south);
\coordinate (railr) at (raill -| sat.south);

% ---- return ----
\node[return box] (return) at ($(railm)+(0,-14mm)$)
    {\textcolor{violet}{\bfseries Return $S$ and $C$}};

% ---- arrows: top pipeline ----
\draw[arr] (input) -- (prep);
\draw[arr] (prep) -- (init);
\draw[arr] (init.south) -- (init.south |- loop.north);

% ---- loop body ----
\draw[arr] (loop.south) -- (learned.north);
\draw[arr] (learned.east) -- node[above] {Yes} (prune.west);
\draw[arr] (learned.south) -- node[right] {No} (solve.north);
\draw[arr, rounded corners=5pt] (solve.south) |- (unsat.east);
\draw[arr, rounded corners=5pt] (solve.south) |- (sat.west);

% ---- bottom rail into return ----
\draw[line width=1pt] (unsat.south) -- (raill) -- (railr) -- (sat.south);
\draw[arr] (railm) -- node[right] {all $HS$ tried} (return);

% ---- feedback: next hitting set ----
\coordinate (fbcol) at ($(sat.east)+(14mm,0)$);
\draw[arr, rounded corners=5pt]
    (prune.east) -- (prune.east -| fbcol) |- (loop.east);
\draw[arr, rounded corners=5pt]
    (railr) -- node[below] {next $HS$} (railr -| fbcol) |- (loop.east);

\end{tikzpicture}
}
    \caption{Factor branching procedure of Algorithm~\ref{alg:Branching}. Refinement produces terminal branches, which the exact solver either represents or certifies as UNSAT.}
    \label{fig:branching_flowchart}
\end{figure}

% Optimizations: found primes absorb later branches
A found branch solution can prune later branches. If this optimization is used, there are two reasons to start with branches that have fewer generators.
First, they are likely to be solved faster. Secondly, they are likely to be higher dimensional and more likely to contain the probably lower dimensional later branches.


% Theoretical analysis of the method
\subsection{Theoretical analysis}
\label{subsec:branching-theoretical-analysis}

\begin{theorem}[One-round branch decomposition]
\label{thm:branch-decomposition}
For every finite $Q\subseteq R\setminus\{0\}$,
\[
    V(Q) = \bigcup_{H \in \mathcal{H}(Q)} V(H).
\]
\end{theorem}

\begin{proof}
Let $a\in V(Q)$. Every $q\in Q$ vanishes at $a$, so some factor in $C_q$
does. The factors in $C(Q)$ that vanish at $a$ therefore form a hitting set,
which contains a minimal one $H$, and $a\in V(H)$. Conversely, let $a\in V(H)$
for some $H\in\mathcal H(Q)$. Every $q\in Q$ has a factor in $H$, which
vanishes at $a$, so $a\in V(Q)$.
\end{proof}

\begin{theorem}[Termination of refinement]
\label{thm:refine-termination}
Algorithm~\ref{alg:Refine} terminates on every finite polynomial system.
\end{theorem}

\begin{proof}
One round replaces each branch $Q$ by its children
$\mathcal H(\widetilde Q)$, where $\widetilde Q=\Call{Interreduce}{Q}$,
independently of the other branches. For a finite set $G$, write
$M(G)=\langle\operatorname{LM}(g):g\in G\rangle$.

We first show that every child $H$ of $Q$ is terminal or satisfies
$M(\widetilde Q)\subsetneq M(\widetilde H)$. Every $q\in\widetilde Q$ is a
multiple of some $h\in H$, so $\operatorname{LM}(h)\mid\operatorname{LM}(q)$
and $M(\widetilde Q)\subseteq M(H)$. Suppose that equality holds, and let
$h\in H$ divide $q\in\widetilde Q$. Since
$\operatorname{LM}(h)\in M(\widetilde Q)$, some $q'\in\widetilde Q$ has
$\operatorname{LM}(q')\mid\operatorname{LM}(h)\mid\operatorname{LM}(q)$.
Since $\widetilde Q$ is interreduced, $q'=q$, so $q/h$ has leading
monomial $1$, and $h=q$ because both are monic. Every
element of $H$ is a factor of an element of $\widetilde Q$, and every element
of $\widetilde Q$ has a factor in $H$, so $H=\widetilde Q$. Then $H$ is
interreduced and consists of monic irreducible polynomials, so
$\widetilde H=H$ and $\mathcal H(H)=\{H\}$, that is, $H$ is terminal.
Otherwise $M(\widetilde Q)\subsetneq M(H)\subseteq M(\widetilde H)$, since
interreduction does not shrink the leading-monomial ideal.

Next, we show that the rounds started from $\{Q\}$ reach a fixed point. Since
$R$ is Noetherian, it has no infinite strictly ascending chain of ideals, so
we may assume this for every branch $H$ with
$M(\widetilde Q)\subsetneq M(\widetilde H)$. A terminal $Q$ is a fixed point.
Otherwise, each of the finitely many children $H$ of $Q$ is terminal or has
$M(\widetilde Q)\subsetneq M(\widetilde H)$, so it reaches a fixed point after
some $k_H$ rounds, and $\{Q\}$ reaches one after $1+\max_Hk_H$ rounds. Since
rounds act on each branch separately, $W=\mathcal H(P)$ reaches a fixed point,
and the loop stops.

Finally, every branch of the final $W$ is terminal. Otherwise, choose a
nonterminal $Q\in W$ whose $M(\widetilde Q)$ is minimal among the nonterminal
branches of $W$. Since $W$ is a fixed point, $Q$ is a child of some $Q'\in W$.
A terminal $Q'$ has only the child $Q'$, so $Q'$ is nonterminal, and
$M(\widetilde Q')\subsetneq M(\widetilde Q)$ contradicts the choice of $Q$.
\end{proof}

\begin{lemma}[Refinement preserves solutions]
\label{lem:recursive-refinement}
Every round of Algorithm~\ref{alg:Refine} preserves
$\bigcup_{Q\in W}V(Q)$. Consequently,
\[
    V(P)=\bigcup_{H\in\mathcal L(P)}V(H)
    \qquad\text{and}\qquad
    \sqrt{\langle P\rangle}=\bigcap_{H\in\mathcal L(P)}\sqrt{\langle H\rangle}.
\]
\end{lemma}

\begin{proof}
Interreduction keeps $V(Q)$, so Theorem~\ref{thm:branch-decomposition} gives
$V(Q)=\bigcup_{H\in\mathcal H(\widetilde Q)}V(H)$ for every branch $Q$, with
$\widetilde Q=\Call{Interreduce}{Q}$, and a round keeps the union. The optional $\Call{Prune}{}$ removes only branches
whose solution sets lie in that of a kept branch. The loop starts from
$W=\mathcal H(P)$, so Theorem~\ref{thm:branch-decomposition} and induction on
the rounds give the first identity. By the Nullstellensatz,
$\sqrt{\langle P\rangle}=I(V(P))$ and $\sqrt{\langle H\rangle}=I(V(H))$, and
the vanishing ideal of a finite union is the intersection of the vanishing
ideals, which gives the second.
\end{proof}

\begin{theorem}[Factor branching]
\label{thm:branching-correctness}
If $\mathcal B$ returns on every branch in $\mathcal L(P)$, then
Algorithm~\ref{alg:Branching} terminates, and
\[
    V(P)=\bigcup_{H\in S}V(H)
    \qquad\text{and}\qquad
    \sqrt{\langle P\rangle}=\bigcap_{H\in S}\sqrt{\langle H\rangle}.
\]
\end{theorem}

\begin{proof}
By Theorem~\ref{thm:refine-termination}, $\Call{Refine}{P}$ returns the finite
set $\mathcal L(P)$, so the algorithm makes finitely many solver calls, each of
which returns. The branches that $\mathcal B$ certifies UNSAT have empty
solution sets, and $\mathcal B$ represents every other branch exactly, so the
first identity follows from Lemma~\ref{lem:recursive-refinement}, and the
second follows from the first by the Nullstellensatz, as in the proof of
Lemma~\ref{lem:recursive-refinement}.
\end{proof}

% Correctness
\begin{theorem}[Prime factor branching]
\label{thm:prime-branching}
If the branch solver $\mathcal B_{\min}$ returns a minimal prime decomposition,
the primes returned by Algorithm~\ref{alg:Branching} that contain no other
returned prime are exactly the minimal primes of $\langle P\rangle$.
\end{theorem}

\begin{proof}
Let $\mathcal M$ be the set of returned primes. The minimal primes of
$\langle H\rangle$ intersect to $\sqrt{\langle H\rangle}$, so
$\bigcap_{\mathfrak p\in\mathcal M}\mathfrak p=\sqrt{\langle P\rangle}$ by
Lemma~\ref{lem:recursive-refinement}. Let $\mathfrak q$ be a minimal prime of
$\langle P\rangle$. It contains the finite intersection
$\bigcap_{\mathfrak p\in\mathcal M}\mathfrak p$, hence some
$\mathfrak p\in\mathcal M$~\cite[Prop.~1.11]{atiyah1969introduction}. Since
$\mathfrak p$ contains $\langle P\rangle$, minimality gives
$\mathfrak p=\mathfrak q$, and $\mathfrak q$ contains no other element of
$\mathcal M$. Conversely, every $\mathfrak p\in\mathcal M$ contains
$\sqrt{\langle P\rangle}$, the intersection of the minimal primes of
$\langle P\rangle$, and hence one of them, which lies in $\mathcal M$. So an
element of $\mathcal M$ that contains no other one is a minimal prime of
$\langle P\rangle$.
\end{proof}
